\section{Why Relay Fails Today: Anatomy of a Coordination Failure}
\label{sec:why-relay-fails}

Relay driving is not new. UPS, FedEx Ground, and BNSF Railway all operate relay-style crew changes within their proprietary driver pools. UPS's Worldport hub in Louisville coordinates relay operations for approximately 200,000 packages per day using a driver workforce organized into fixed shift assignments. The difference between UPS relay and the missing inter-company relay marketplace is identical to the difference between a private corporate jet program and a commercial airline: one solves the coordination problem internally through vertical integration; the other solves it through a market with access rules, pricing, and neutral platform infrastructure.

The failure of inter-company relay to emerge spontaneously over the 60 years since the interstate highway system was completed requires explanation. Four structural barriers operate simultaneously and reinforce each other.

\subsection{Carrier Fragmentation and the Driver Pool Problem}

The US trucking industry is structurally inhospitable to relay. The FHWA carrier census counts more than 70,000 registered for-hire motor carriers; 90\% operate fewer than 10 power units \citep{FHWA_freight2023}. The median carrier fleet size is 3--4 trucks. A relay operation requires that a carrier have a driver available at the relay hub when the truck arrives --- not eventually, but within a 20-minute window that prevents trailer detention charges from accumulating. For a 4-truck carrier, the probability that a driver happens to be on-shift and available at the exact hub at the exact time is approximately zero. The carrier would need to station a driver at the hub speculatively, at a cost of \$175--250/day in wages and benefits \citep{BLS_trucking2023}, on the chance that one of its 4 trucks would arrive during that driver's shift. The economics of speculative driver positioning do not work for small carriers.

Large carriers solve this problem through scale: J.B. Hunt, Werner, and Amazon Freight each operate fleets exceeding 10,000 tractors and can statistically guarantee hub presence \citep{FHWA_freight2023}. But large-carrier relay solves only large-carrier freight; the 90\% of the market operated by small carriers is excluded. The relay marketplace solves this problem structurally: hub operators pool relay drivers across all carriers, so a 4-truck carrier gets hub-presence at the same quality of service as J.B. Hunt. This is precisely what an airport ground handler provides: a 12-seat regional airline operating 2 flights per day at an airport gets the same baggage handling quality as United Airlines, because the ground handler's driver pool is shared infrastructure.

\subsection{The Regulatory Gap in 49 CFR Part 395}

Hours of Service (HOS) regulation under 49 CFR Part 395 governs driver duty periods, rest requirements, and electronic logging device (ELD) recording standards \citep{FMCSA_HOS2022}. The regulation was written assuming one driver per vehicle per trip. It does not define a ``relay terminal'' as a location where driver responsibility transfers between ELD systems. It does not specify a streamlined pre-trip inspection protocol for a truck that was inspected by a departing driver 6 hours prior. It does not define an ``inter-carrier HOS transfer'' that would allow Driver B's ELD to pick up the load documentation created by Driver A's ELD at a different carrier's USDOT number.

This is a regulatory gap, not a prohibition. Nothing in 49 CFR Part 395 explicitly prohibits relay; the regulation simply does not contemplate it. The practical consequence is that every relay operation must be structured as two separate trips (Driver A delivers the trailer to the hub; hub becomes the shipper; Driver B picks up as a new load) with full pre-trip documentation, insurance re-certification, and bill of lading transfer at each handoff. The overhead of this workaround --- approximately 45--90 minutes of administrative time per handoff --- eliminates the time-efficiency benefit that makes relay valuable in the first place.

The required FMCSA rulemaking is modest in scope. It needs to define three concepts: (1) an approved relay terminal as a location where HOS-compliant driver transfer occurs; (2) a relay pre-trip inspection as a condensed 5-item inspection (brakes, lights, tires, coupling, mirrors) where the vehicle health has been pre-validated by OBD diagnostics within the prior 8 hours; and (3) an inter-carrier ELD handoff as an electronic transaction in which Driver B's ELD inherits the load identification, DOT numbers, and HOS-to-date data from Driver A's ELD via a platform API. None of these require statutory change. All are achievable via rulemaking under FMCSA's existing authority \citep{FMCSA_HOS2022}. The standard rulemaking timeline --- NPRM, comment period, final rule --- is 18--24 months.

\subsection{The First-Mover Disadvantage}

Relay hub infrastructure requires capital investment before the carrier network that would use it exists. A hub operator building a 750-dock relay facility in Atlanta invests \$5 million in infrastructure (land, fuel, driver facilities, dispatch technology, and truck monitoring systems) before a single carrier commits to routing its trucks through the hub. The carriers, in turn, will not commit to hub routing until they know the hub will be reliably staffed --- which requires the hub operator to have completed the infrastructure investment. This is a classic bilateral hold-up problem \citep{Williamson1979}, and it has suppressed relay hub investment for the same reason that chicken-and-egg platform problems suppressed credit card networks before Visa and Mastercard solved them through simultaneous market-side subsidization.

The hold-up is compounded by the absence of a neutral coordination mechanism. In the airline industry, the FAA designates airport facilities as public-use infrastructure, removing the first-mover risk for ground handlers. In railroads, trackage rights agreements create the legal basis for one railroad to commit loads to infrastructure owned by another. In trucking, no equivalent designation or legal framework exists for relay hub infrastructure. The government role needed is narrow: under the Infrastructure Investment and Jobs Act (IIJA), FHWA has existing authority to designate portions of the National Highway Freight Network as National Freight Relay Zones (NFRZ) --- facilities eligible for priority funding and operational coordination status. An NFRZ designation at T1 hub locations does not build the hub; it removes the first-mover risk by signaling durable government commitment to the hub's operational role.

\subsection{The Hidden Cost Illusion}

The most persistent barrier to relay adoption is a systematic accounting error in carrier cost models. When a carrier evaluates a solo long-haul driver versus relay, it typically computes:

\begin{align*}
C_{\text{solo}} &= r_{\text{mile}} \times d = \$0.52/\text{mi} \times 2{,}800\text{ mi} = \$1{,}456/\text{trip} \\
C_{\text{relay}} &= n_{\text{drivers}} \times h_{\text{shift}} \times w_{\text{hour}} = 6 \times 7\text{h} \times \$25/\text{h} = \$1{,}050/\text{trip}
\end{align*}

The relay savings of \$406 per trip appear real, but the carrier also sees an estimated coordination cost of \$50--150 per handoff (scheduling, documentation, hub fee) across 5 handoffs per NY--LA trip, eroding the apparent advantage. The visible conclusion: relay saves \$156--\$656 per trip depending on coordination efficiency --- attractive but not transformative.

What the visible calculation omits is asset utilization. A solo driver under HOS regulations drives 11 hours, sleeps 10 hours, and completes the cycle in 21 hours. The truck is idle 10 out of 21 hours --- a 48\% utilization rate \citep{FMCSA_HOS2022}. A relay-operated truck runs continuously; the utilization rate approaches 98\% (accounting for fueling and inspection stops). At \$180,000 per Class 8 tractor, the capital implication is significant: a carrier delivering the same freight-ton-miles via relay needs approximately half the truck fleet it needs under solo operations. On a 10-truck fleet operating 260 days per year, this represents \$900,000 in truck capital that is not deployed --- the equivalent of 5 trucks freed up for other routes or sold.

The full cost comparison, incorporating asset utilization, is:

\begin{align*}
\Delta C_{\text{truck capital}} &= (1 - 0.48/0.98) \times \$180{,}000 \times r_{\text{capital}} \approx \$8{,}700/\text{truck/year}
\end{align*}

where $r_{\text{capital}}$ is the annualized capital recovery factor at 7\% over 5 years (approximately 0.243). Spread across the truck's annual route volume, the capital savings per trip exceed the visible relay coordination cost premium on any route longer than 1,500 miles. Additionally, long-haul solo drivers command signing bonuses of \$5,000--\$15,000 in the current shortage market \citep{ATRI_shortage2024}, whereas regional relay drivers --- who work predictable shifts and sleep at home --- face no comparable shortage premium. The per-driver hiring and retention cost advantage of relay over solo long-haul is not captured in any standard carrier cost model.

The accounting error is systematic because carrier cost systems were built for solo operations. Relay requires a different cost accounting structure: one that capitalizes truck fleet size as a function of utilization rate rather than route volume. No carrier has built this accounting, because no relay marketplace has existed to make it necessary.

\subsection{Summary of Barriers}

Table~\ref{tab:barriers} summarizes the four barriers and the corresponding platform intervention required for each.

\begin{table}[ht]
\centering
\caption{Coordination failure barriers and relay marketplace interventions}
\label{tab:barriers}
\begin{tabular}{p{3cm}p{5cm}p{5cm}}
\toprule
\textbf{Barrier} & \textbf{Mechanism} & \textbf{Intervention} \\
\midrule
Carrier fragmentation & 70k+ carriers cannot individually sustain hub driver pools & Hub operator pools drivers across all carriers; any carrier pays per-swap \\
Regulatory gap & 49 CFR Part 395 does not define relay terminal, pre-trip, or ELD transfer & FMCSA rulemaking: 3 definitions, 18--24 months \\
First-mover disadvantage & Hub investor and carrier network must commit simultaneously & NFRZ designation under IIJA removes investment risk \\
Hidden cost illusion & Carrier cost models omit asset utilization savings from relay & Industry education + relay marketplace publishes utilization analytics \\
\bottomrule
\end{tabular}
\end{table}
